% !TeX program = xelatex
% ============================================================================
%  第 13 课　它有什么用　—— 教师版讲义
%  与 02-课件/第13课-它有什么用/第13课-它有什么用.tex 同步
% ============================================================================

\jhlesson{13}{它有什么用}{时长 45 分钟\quad·\quad 教具：无\quad·\quad 本课不发单页}

\begin{mubiao}
  \begin{itemize}
    \item 会用费马小定理快速求大幂的余数。
    \item 能回答第 1 课留下的问题：\(999999\) 为什么能被 \(7\) 整除。
    \item 知道身份证校验码怎么算，RSA 大致在做什么。
  \end{itemize}
\end{mubiao}

\noindent{\small 本课节奏：0—6 回顾｜6—16 回答第 1 课的问题｜16—28 快速求幂｜28—38 身份证校验码（含「用你自己的号码验一遍」，约 3 分钟）｜38—45 RSA（不够就按下文预案压到 3 分钟）。}

\section*{一、回顾（0—6 分钟）}

把费马小定理原文再读一遍：\(p\) 是质数、\(a\) 不是 \(p\) 的倍数时，\(a^{\,p-1}\equiv1\pmod p\)。

一句话过渡：工具磨好了，今天先拿它去解决一件旧事。

\section*{二、回答第 1 课的问题（6—16 分钟）}

\begin{caozuo}
  指着黑板最右边那行「\(999999\div7=?\)　第 13 课回来回答」。
  先别说答案，一步一步带全班写完，写完再走到黑板右边，当场把这行字划掉，写上「已解决」。
\end{caozuo}

第一步，改写：
\[
  999999=1000000-1=10^{6}-1 .
\]

第二步，用定理：\(7\) 是质数，\(10\) 不是 \(7\) 的倍数，取 \(p=7\)、\(a=10\)：
\[
  10^{\,7-1}=10^{6}\equiv1\pmod 7 .
\]

第三步，代回去：
\[
  10^{6}-1\equiv1-1=0\pmod 7 .
\]

所以 \(7\) 整除 \(999999\)。整门课埋的那个问号，在这里收口。

\begin{zhu}
  这里要停两秒，让学生自己体会一下：我们从头到尾没有做除法，也没有用计算器。
\end{zhu}

\section*{三、快速求幂（16—28 分钟）}

\begin{caozuo}
  求 \(2^{100}\) 除以 \(7\)，一步一步写在黑板上：
  \[
    100=6\times16+4,\qquad
    2^{100}=(2^{6})^{16}\times2^{4}\equiv1^{16}\times16\equiv2\pmod 7 .
  \]
  每一步都让全班用手指报下一个结果。
\end{caozuo}

再练一道：求 \(3^{200}\) 除以 \(11\)。

\[
  3^{10}\equiv1\pmod{11},\qquad
  200=10\times20,\qquad 3^{200}=(3^{10})^{20}\equiv1\pmod{11}.
\]

\begin{zhu}
  做题的套路只有一句：\emph{把指数除以 \(p-1\) 看余数}。
  因为 \(a^{\,p-1}\equiv1\)，多出来的整圈可以直接扔掉。
\end{zhu}

\begin{caozuo}
  手指回答：\(2^{100}\) 除以 \(7\) 余几？（\(2\)）
  求 \(5^{2026}\) 除以 \(7\)，第一步看什么？（先算 \(2026\div6\) 的余数；\(2026=6\times337+4\)，
  所以答案与 \(5^{4}\equiv2\) 相同）
\end{caozuo}

\begin{zhu}
  时间弹性：若到 \(34\) 分钟还没讲完快速求幂，\(3^{200}\) 这一题让学生课后做——
  \textbf{身份证一节不能砍}，那是学生最想看的一段。
\end{zhu}

\section*{四、身份证校验码（28—38 分钟）}

身份证前 \(17\) 位，每一位乘一个固定的权重，加起来，再除以 \(11\) 取余；余数查表换成最后一位。

\begin{center}
\begin{tabular}{@{}lc*{17}{c}@{}}
  \toprule
  位数 & \(1\) & \(2\) & \(3\) & \(4\) & \(5\) & \(6\) & \(7\) & \(8\) & \(9\) & \(10\) & \(11\) & \(12\) & \(13\) & \(14\) & \(15\) & \(16\) & \(17\) \\
  \midrule
  权重 & \(7\) & \(9\) & \(10\) & \(5\) & \(8\) & \(4\) & \(2\) & \(1\) & \(6\) & \(3\) & \(7\) & \(9\) & \(10\) & \(5\) & \(8\) & \(4\) & \(2\) \\
  \bottomrule
\end{tabular}
\end{center}

\begin{center}
\begin{tabular}{@{}l*{11}{c}@{}}
  \toprule
  余数 & \(0\) & \(1\) & \(2\) & \(3\) & \(4\) & \(5\) & \(6\) & \(7\) & \(8\) & \(9\) & \(10\) \\
  \midrule
  校验码 & \(1\) & \(0\) & X & \(9\) & \(8\) & \(7\) & \(6\) & \(5\) & \(4\) & \(3\) & \(2\) \\
  \bottomrule
\end{tabular}
\end{center}

\begin{caozuo}
  取一个虚构号码的前 \(17\) 位 \(3\,5\,0\,6\,0\,4\,2\,0\,1\,3\,0\,5\,2\,0\,1\,2\,3\)，当着全班一步一步算：
  \[
    21+45+0+30+0+16+4+0+6+9+0+45+20+0+8+8+6=218,
  \]
  \[
    218=11\times19+9 \;\Longrightarrow\; \text{余数 }9 \;\Longrightarrow\; \text{校验码 }3 .
  \]
  所以完整的编号是 \(350604201305201233\)。
\end{caozuo}

\begin{zhu}
  板书时一位一位对齐着写，不要跳步。学生最想看的就是「数字怎么变成最后那一位」。
\end{zhu}

\begin{zhu}
  这个号码是编出来的。前六位 \(350604\) 是漳州市龙海区的行政区划代码——
  漳州港开发区就在龙海区境内，所以拿它当例子，学生一眼认得出来。
  （龙海撤市设区以前这个代码是 \(350681\)，也不必对学生讲。）
\end{zhu}

\begin{caozuo}
  \textbf{新增环节：让学生用自己带来的号码验一遍（学生版同题，约 \(3\) 分钟）。}
  \begin{enumerate}
    \item 先请学生把号码\textbf{前 \(17\) 位}抄在讲义那张表里。
      \textbf{提醒：只写号码，不要写名字}——讲义会收走、会传阅；
      不想写自己号码的，就用刚才那个虚构号码。
    \item 要求\textbf{分三段填小计}：第 \(1\)—\(6\) 位（地址码）、第 \(7\)—\(14\) 位（出生日期）、
      第 \(15\)—\(17\) 位（顺序码）。三段分别加、最后再合，能少一半算错。
    \item 走一圈看：算得快的同学，请他和同桌对一对；算不完的\textbf{不要催}，
      让他先把第 \(3\) 段（第 \(15\)—\(17\) 位，三项）填上，其余课后补完。
    \item 收尾问一句：「算出来的最后一位，和你号码上印的最后一位一样吗？」
      一样就说明这张号码是「对」的——校验码就是干这个用的。
  \end{enumerate}
  \textbf{若时间紧，只做最小版本「验第 \(17\) 位」：}让学生算出 \(d_{17}\times2\) 这一个乘积，
  告诉他「这只是 \(17\) 项里的一项，其余 \(16\) 项用同一张表课后补完，就能验出你自己的号码」。
\end{caozuo}

\begin{zhu}
  \textbf{答案口径（示例号 \(350604201305201233\)）：}三段小计分别是
  第 \(1\)—\(6\) 位 \(=112\)、第 \(7\)—\(14\) 位 \(=84\)、第 \(15\)—\(17\) 位 \(=22\)，合计 \(218\)。
  学生自己的号码没有统一答案，但\textbf{只要前六位是 \(350604\)，第一段小计一定是 \(112\)}——
  可以在黑板上先把这个数写死，学生就只剩两段要算。
  学生报出来的总和若对不上，多半是抄错位或是三段加错：让他先检查
  「前六位是不是 \(350604\)」「最后一位抄对了没有」。
\end{zhu}

\begin{zhu}
  \textbf{⚠️ 上课前先扫一眼：班上如果混着别的地区的学生，"\(112\)"这个数对他们不成立。}
  第 \(1\)—\(6\) 位是地址码，只有前六位恰好是 \(350604\)（漳州龙海）的同学，第一段小计才是 \(112\)；
  换成别的代码这个数就变了（例如 \(350681\to164\)、\(350602\to104\)）。
  \textbf{做法}：先问一句"有没有同学的前六位不是 \(350604\)？"，
  \begin{itemize}
    \item 全班基本都在龙海 \(\to\) 把 \(112\) 写死，按上面的流程走；
    \item 有零星几个外地同学 \(\to\) 让他们\textbf{自己算第一段}（也就是多算 6 项，不影响体验），
      或者干脆请他们用虚构号码；
    \item 不确定 \(\to\) 就\textbf{不给 \(112\)}，让所有人自己算第一段，"分三段填小计"这个脚手架已经够用了。
  \end{itemize}
  \textbf{一句话}：\(112\) 是省事用的，不是必须的；宁可让学生多算六项，也不要让他觉得"我的号码算不出 \(112\)，是不是我错了"。
\end{zhu}

\begin{zhu}
  \textbf{课后选做「只算 \(9\) 项」：}让他们用示例号只算「前 \(6\) 位 \(+\) 后 \(3\) 位」共 \(9\) 项
  （\(112+22=134\)），再用已知的总和 \(218\) 倒推中间 \(8\) 项：\(218-134=84\)。
  这样既少算了 \(8\) 项，又验证了「分段相加」这件事。
\end{zhu}

追问一句：为什么这一位能防错？——改掉前 \(17\) 位里任意一位，加权和几乎一定会变，
余数变了，校验码就对不上。

\section*{五、RSA 的故事（38—45 分钟）}

\noindent 这一节只讲概览，不写公式，也不带学生算。

\begin{caozuo}
  用生活化的比喻开场：想让别人把秘密寄给你，又不想把钥匙交给别人，怎么办？
  于是造一把「谁都能拿到的锁」（公钥）和一柄「只有你有的钥匙」（私钥）：谁都锁得上，只有你打得开。
\end{caozuo}

\begin{zhu}
  为什么别人破解不了？因为钥匙是用两个很大的质数\emph{乘}出来的，
  而想把一个很大的数\emph{拆回}两个质数，非常困难——「合起来容易，拆开难」。
  背后是费马小定理的推广（欧拉定理）。本课只讲结论，不证。
\end{zhu}

\begin{zhu}
  数学史（务必说准）：
  \begin{itemize}
    \item \(1977\) 年，美国麻省理工学院的三位学者 Rivest、Shamir、Adleman 发表这套方法，
          取名 \textbf{RSA}（三人姓氏的首字母）；\(2002\) 年三人获得图灵奖；
    \item 更早的 \(1973\) 年，英国情报机构 GCHQ 的 Clifford Cocks 也独立想出过同样的方法，
          因列为机密直到 \(1997\) 年才公开——所以公开的第一份仍算 RSA；
    \item 网页、手机、银行卡上用到的加密，大多是 RSA 这一类方法（或它的近亲）。
  \end{itemize}
\end{zhu}

\begin{zhu}
  时间不够时，本节可以压到 3 分钟：只讲比喻 + 一句「拆开难、合起来易」+ 一句数学史（\(1977\) 年、三人、图灵奖）。
  时间富余时，再加一句「规模」：一个很小的例子 \(33=3\times11\) 一眼就能拆开；
  真正在用的 \(n\) 有几百位长，拆开它，今天的计算机也算不完。
\end{zhu}

收尾用一句话：手机里、网页上，每一次 https 用的都是这类\textbf{绕圈的算术}。

\begin{xiaojie}
  \begin{itemize}
    \item \(999999=10^{6}-1\)，由 \(10^{6}\equiv1\pmod7\) 得 \(7\) 整除 \(999999\)。
    \item 快速求幂：先算指数除以 \(p-1\) 的余数，\(2^{100}\equiv2\pmod7\)、\(3^{200}\equiv1\pmod{11}\)。
    \item 身份证校验码、RSA 加密，用的都是「取余」这套算术。
  \end{itemize}
\end{xiaojie}

\noindent\textbf{下节课}：复习、总结与展望——把整门课串成一条线。